\documentclass[11pt]{article}
\usepackage[margin=0.9in]{geometry}
\usepackage{amsmath,amssymb,booktabs}
\usepackage[hidelinks]{hyperref}
\title{Standard Model quantum numbers from the $E_6\times SU(3)$ chirality pair}
\author{Generated by \texttt{lisi/make\_report.py} (sympy $1.14.0$)}
\date{}
\begin{document}
\maketitle

\begin{abstract}
Starting only from the 240 roots of $E_8$, we recover the branching
$27\to 1+16+10$ under $SO(10)\times U(1)$, the Standard Model content of the
$16$, the hypercharge generator, and the vanishing of the linear and cubic
$U(1)$ sums.  Every statement is an exact rational computation; see the
scripts in \texttt{lisi/}.  Nothing here is a derivation of the physical
Standard Model, only a check that the chirality-theorem pair contains it.
\end{abstract}

\section{Setup}
Roots of $E_8$ are stored in doubled integer coordinates.  The horizontal
$SU(3)_M$ is the $A_2$ spanned by two roots with inner product $-1$, and
$E_6$ is its centralizer.  The $240$ roots split as $72+6+6\cdot27$, the root
form of
\[
  248=(78,1)+(1,8)+(27,3)+(\overline{27},\bar 3).
\]

\section{From $E_6$ to $SO(10)$}
Fix one weight $w_0$ of the $27$ (in true coordinates, $\langle w_0,w_0\rangle=\frac{4}{3}$):
\[
  w_0=\left[\begin{matrix}- \frac{2}{3} & - \frac{2}{3} & - \frac{2}{3} & 0 & 0 & 0 & 0 & 0\end{matrix}\right] .
\]
The $E_6$ roots orthogonal to $w_0$ form a system of $40$ roots, that of
$SO(10)$, and the $27$ splits by $\langle w,w_0\rangle$:
\[
  27 = \underbrace{\mathbf{1}}_{4/3}\;+\;\underbrace{\mathbf{16}}_{1/3}\;+\;\underbrace{\mathbf{10}}_{-2/3}.
\]

\section{Standard Model content}
Inside $SO(10)$ we select an $A_2\times A_1$ subsystem (roots $a,b$ and $g$)
for which the $16$ branches as one Standard Model generation:
\[
  a=\left[\begin{matrix}0 & 0 & 0 & 1 & 1 & 0 & 0 & 0\end{matrix}\right],\quad b=\left[\begin{matrix}0 & 0 & 0 & -1 & 0 & 1 & 0 & 0\end{matrix}\right],\quad g=\left[\begin{matrix}0 & 0 & 0 & 0 & 0 & 0 & 1 & 1\end{matrix}\right] .
\]
The hypercharge generator is the vector orthogonal to $A_2\times A_1$, to
$w_0$ and to the horizontal $A_2$ that reproduces the Standard Model
values; a weight $w$ has $Y=\langle w,Y\rangle$ and $T_3=\langle w,g\rangle/2$ with
\[
  Y=\left[\begin{matrix}0 & 0 & 0 & \frac{1}{3} & - \frac{1}{3} & \frac{1}{3} & \frac{1}{2} & - \frac{1}{2}\end{matrix}\right] .
\]
\subsection*{The $\mathbf{16}$}
\begin{center}
\begin{tabular}{llccr}
\toprule
field & $\mathrm{SU}(3)_c\times\mathrm{SU}(2)_L$ & $Y$ & $Q=T_3+Y$ & weights \\
\midrule
$Q_L$ & $(\mathbf{3},\mathbf{2})$ & $\frac{1}{6}$ & $- \frac{1}{3},\ \frac{2}{3}$ & 6 \\
$u^c$ & $(\overline{\mathbf{3}},\mathbf{1})$ & $- \frac{2}{3}$ & $- \frac{2}{3}$ & 3 \\
$d^c$ & $(\overline{\mathbf{3}},\mathbf{1})$ & $\frac{1}{3}$ & $\frac{1}{3}$ & 3 \\
$L$ & $(\mathbf{1},\mathbf{2})$ & $- \frac{1}{2}$ & $-1,\ 0$ & 2 \\
$\nu ^c$ & $(\mathbf{1},\mathbf{1})$ & $0$ & $0$ & 1 \\
$e^c$ & $(\mathbf{1},\mathbf{1})$ & $1$ & $1$ & 1 \\
\bottomrule
\end{tabular}
\end{center}
\subsection*{The $\mathbf{10}$ and the singlet}
\begin{center}
\begin{tabular}{llccr}
\toprule
field & $\mathrm{SU}(3)_c\times\mathrm{SU}(2)_L$ & $Y$ & $Q=T_3+Y$ & weights \\
\midrule
$D$ & $(\mathbf{3},\mathbf{1})$ & $- \frac{1}{3}$ & $- \frac{1}{3}$ & 3 \\
$D^c$ & $(\overline{\mathbf{3}},\mathbf{1})$ & $\frac{1}{3}$ & $\frac{1}{3}$ & 3 \\
$H_d$ & $(\mathbf{1},\mathbf{2})$ & $- \frac{1}{2}$ & $-1,\ 0$ & 2 \\
$H_u$ & $(\mathbf{1},\mathbf{2})$ & $\frac{1}{2}$ & $0,\ 1$ & 2 \\
\bottomrule
\end{tabular}
\end{center}
The singlet has $Y=0$ and $Q=0$.

\paragraph{Cross-check.}  On the $40$ roots of $SO(10)$ this $Y$ takes the
values of the $SU(5)$ adjoint pattern: eight roots at $0$, twelve at $\pm5/6$,
twelve at $\pm1/6$, six at $\pm2/3$ and two at $\pm1$.  This is the standard
hypercharge, not a flipped one.

\section{Anomaly sums}
Exact sums over the $\mathbf{16}$:
\begin{align*}
  \sum Y &= 3 \left(- \frac{2}{3}\right) + 2 \left(- \frac{1}{2}\right) + 1 \cdot 0 + 6 \cdot \frac{1}{6} + 3 \cdot \frac{1}{3} + 1 \cdot 1 = 0,\\
  \sum Y^3 &= 3 \left(- \frac{8}{27}\right) + 2 \left(- \frac{1}{8}\right) + 1 \cdot 0 + 6 \cdot \frac{1}{216} + 3 \cdot \frac{1}{27} + 1 \cdot 1 = 0 .
\end{align*}
The $\mathbf{10}$ is vector-like, so its sums vanish termwise; the $\mathbf{1}$
carries $Y=0$.  The sums $\sum Q$ and $\sum Q^3$ over the $16$ also vanish, as
the scripts verify.  This is the Standard Model gauge-anomaly cancellation
within a single generation, visible here as a property of the $E_6$ weights.

\section{Three families and triality}
In the chirality theorem the representation is $V=27\boxtimes 3$, so the net
chiral content is three copies of $16$, i.e.\ $48$ Weyl fermions, together with
three vector-like $10$'s and three singlets.  A separate check places the second
$D_4$ of $D_4\times D_4\subset E_8$ inside $E_6$, giving
$27=8_v+8_s+8_c+3\cdot\mathbf 1$, the triality triple used by Lisi for
generations.  Whether the horizontal $3$ and the triality $3$ are related is an
open question that these computations do not settle.

\section{Reproducing}
\begin{verbatim}
python3 lisi/run_all.py          # all checks
python3 lisi/make_report.py      # regenerate this PDF
\end{verbatim}
\end{document}
